% ============================================================
%  preamble.tex — Package & Setting untuk Skripsi NusaNara
% ============================================================

% --- Encoding & Bahasa ---
\usepackage[T1]{fontenc}
\usepackage[utf8]{inputenc}
\usepackage[bahasa]{babel}

% --- Font ---
\usepackage{lmodern}

% --- Layout & Margin ---
\usepackage[a4paper, top=3cm, bottom=3cm, left=4cm, right=3cm]{geometry}
\usepackage{setspace}
\onehalfspacing              % Spasi 1.5 (standar skripsi)

% --- Gambar & Grafik ---
\usepackage{graphicx}

\usepackage{float}           % [H] placement
\graphicspath{{figures/}}    % Folder gambar default

% --- Tabel ---
\usepackage{booktabs}        % Garis tabel profesional
\usepackage{tabularx}        % Tabel lebar fleksibel
\usepackage{longtable}       % Tabel multi-halaman
\usepackage{multirow}        % Sel gabungan
\usepackage{array}
\usepackage{xcolor}
\usepackage{colortbl}        % Warna sel tabel

% --- Matematika & Rumus ---
\usepackage{amsmath}
\usepackage{amssymb}
\usepackage{amsfonts}
\usepackage{bm}              % Bold math

% --- Kode Program ---
\usepackage{listings}
\usepackage{textcomp}        % Supaya backtick (`) dirender sebagai backtick asli (upquote)
\usepackage{xcolor}
\definecolor{codebg}{RGB}{245,245,245}
\definecolor{codecomment}{RGB}{100,100,100}
\definecolor{codestring}{RGB}{163,21,21}
\definecolor{codekeyword}{RGB}{0,0,255}


\lstdefinelanguage{JavaScript}{
  keywords={typeof, new, true, false, catch, function, return, null, catch, switch, var, if, in, while, do, else, case, break, const, let},
  keywordstyle=\color{blue}\bfseries,
  sensitive=false,
  comment=[l]{//},
  morecomment=[s]{/*}{*/},
  morestring=[b]',
  morestring=[b]"
}

\lstset{
  upquote=true,              % Memastikan kutipan dan backtick dirender lurus
  backgroundcolor=\color{codebg},
  basicstyle=\small\ttfamily,
  commentstyle=\color{codecomment}\itshape,
  stringstyle=\color{codestring},
  keywordstyle=\color{codekeyword}\bfseries,
  breaklines=true,
  frame=single,
  rulecolor=\color{gray!40},
  numbers=left,
  numberstyle=\tiny\color{gray},
  numbersep=5pt,
  showstringspaces=false,
  tabsize=2,
  captionpos=b
}

% --- Hyperlink & PDF Metadata ---
\usepackage[hidelinks]{hyperref}
\hypersetup{
  pdftitle={Revisi Skripsi NusaNara - Bab 2, 3, 4},
  pdfauthor={Muhammad Mirza Kurniawan},
  pdfsubject={Sistem Bimbingan Karier Berbasis RAG dan LLM},
}

% --- Caption & Referensi ---
\usepackage{caption}
\captionsetup{
  font=small,
  labelfont=bf,
  format=plain,
  justification=centering
}

% --- Penomoran ---
\usepackage{enumitem}        % List yang fleksibel

% --- Chapter & Section Heading ---
\usepackage{titlesec}
\titleformat{\chapter}[display]
  {\normalfont\large\bfseries\centering}
  {BAB \thechapter}{0pt}{\large}
\titlespacing{\chapter}{0pt}{-20pt}{20pt}

% --- Paragraf ---
\setlength{\parindent}{1.25cm}
\setlength{\parskip}{0pt}

% --- Footer/Header ---
\usepackage{fancyhdr}
\pagestyle{fancy}
\fancyhf{}
\fancyhead[R]{\thepage}
\fancyhead[L]{\leftmark}
\renewcommand{\headrulewidth}{0.4pt}

% --- Utility ---
\usepackage{lipsum}

\newcommand{\missingfig}[1]{\begin{center}\fbox{\parbox{0.8\linewidth}{\centering\small\textbf{[GAMBAR PLACEHOLDER]}\\\textit{#1}\\\textit{Sisipkan dari dokumen asli}}}\end{center}}          % Dummy text (hapus di versi final)
\usepackage{csquotes}
